\chapter{Discusión}
\label{ch:discusion}

Los resultados sustentan una afirmación deliberadamente acotada. Este capítulo separa
lo que el marco demuestra con solidez de lo que demuestra solo parcialmente, y
argumenta que la distinción depende de dos preguntas distintas que el trabajo previo a
menudo confunde. La tesis central es esta: el marco aporta un \emph{diagnóstico de
brecha de demanda validado} y un \emph{generador de corredores caracterizado y
parcialmente exitoso}, y estas dos contribuciones tienen pesos probatorios distintos.

\section{Dos preguntas, dos capas}
\label{sec:disc-twoquestions}

Conviene distinguir dos preguntas que se le podrían plantear a un marco de planeación
orientado por datos. Pregunta~A: ¿reproduce el generador \emph{corredores que
realmente se construyeron}? Pregunta~B: ¿produce el generador \emph{buenos
corredores}---líneas que atienden demanda real y no atendida de forma eficiente y sin
redundancia? Buena parte de la literatura sobre ubicación de estaciones con aprendizaje
automático responde implícitamente una versión de la Pregunta~A, tratando la
coincidencia con la red existente como validación. Pero A es un proxy débil y
asimétrico de B, porque las líneas construidas no están garantizadas como óptimas: se
eligen bajo restricciones políticas, financieras y de disponibilidad de suelo que el
marco no modela. La coincidencia, por tanto, corrobora; el desacuerdo es evidencia
débil de error. Las dos capas de la metodología corresponden a estas preguntas ---la
capa de diagnóstico (W1--W4) se valida contra evidencia de tipo B, mientras que el
generador (W5--W6) se prueba contra ambas--- y tienen éxito en grados distintos.

\section{La capa de diagnóstico es la contribución sólida}
\label{sec:disc-diagnostic}

El diagnóstico de brecha de cobertura (\cref{ch:metodologia}, \cref{sec:w3}) es el
resultado más firme del marco. Su contribución metodológica es romper la circularidad
que socava los objetivos de ``tiene parada'': al definir la variable dependiente como
la razón entre la demanda estimada de forma independiente y la oferta medida por
\textsc{gtfs} (\cref{eq:gap}), el modelo aprende dónde se \emph{necesita} el transporte
en lugar de dónde \emph{existe}. El reentrenamiento supervisado confirma que este
objetivo es aprendible solo a partir de características del lugar, sin fuga de
información de la oferta (\cref{tab:retrain}).

Tres líneas de evidencia sustentan el diagnóstico. Primera, y más importante, se
corrobora \emph{fuera de muestra}: la instantánea \textsc{gtfs} de 2024 es anterior a
la Línea~4 de SITEUR, por lo que el corredor se trata como desatendido, y sin embargo
sus \ageb{}s muestran \num{1.6} veces la tasa metropolitana de brecha alta
(\cref{sec:res-w8}). Un corredor que los planificadores profesionales decidieron
construir de forma independiente fue señalado como de alta prioridad por el modelo que
nunca lo vio. Segunda, el diagnóstico se \emph{transfiere}: produce una lectura
plausible y diferenciada de la demanda no atendida en dos ciudades adicionales
(\cref{ch:transferibilidad}), ordenándolas con sensatez por proporción de brecha alta.
Tercera, es \emph{interpretable}: la atribución SHAP identifica zonas densas, de alta
necesidad y ricas en empleo pero desatendidas, coincidiendo con el patrón geográfico
de la \cref{fig:coverage-gap}. Esta capa es el núcleo defendible y citable de la tesis.

\section{La capa generativa es un éxito parcial y caracterizado}
\label{sec:disc-generative}

El generador de corredores es una contribución más débil pero genuina, y su historia
es instructiva. En su primera forma el generador era esencialmente negativo: de sus
corredores factibles, solo un tramo corto y degenerado pasaba las pruebas de mérito, y
la factibilidad estaba confundida con la escasez de anclas ---el filtro de
factibilidad seleccionaba por \emph{pocas anclas casi colineales} en lugar de por
buenos corredores. Esto se rastreó a dos defectos acoplados: un modelador que aplanaba
un árbol de expansión mínima en una línea con saltos fantasma, y una métrica de rodeo
medida contra la distancia entre extremos, que penaliza en exceso a un corredor de
cobertura de demanda que legítimamente se curva.

Rearquitecturar el generador en torno a anclas de frontera, un tronco de
\emph{diámetro} del árbol de expansión mínima ensamblado de calles reales y una
compuerta de factibilidad por \emph{rectitud de anclas} (\cref{sec:w6}) cambió el
veredicto. En las pruebas de mérito de tres ejes, el generador ahora produce al menos un corredor
sustantivo, factible y meritorio: \code{W6\_G02} atiende \SI{12.1}{\kilo\metre} y
\num{25} \ageb{}s, de los cuales el \SI{56}{\percent} son de brecha alta (frente a la
línea base de \SI{20.7}{\percent}), no es redundante con la red existente (mejor
Jaccard \num{0.18}) y se ubica en el percentil \num{73} de demanda por kilómetro. Esta
es una línea real, no un tramo.

El éxito es parcial, y de una manera más específica de lo que caracterizaban versiones
anteriores de este capítulo: el veredicto de las pruebas de mérito y la propia función
multiobjetivo W5 de la metodología discrepan sobre cuál corredor es el mejor. La
\cref{tab:pareto-w5} muestra que bajo dominancia estricta de Pareto ---no sensible a
ninguna elección de pesos objetivo--- \code{W6\_G03}, no \code{W6\_G02}, es el único
corredor factible no dominado; \code{W6\_G02} está dominado en los tres objetivos tanto
por \code{W6\_G00} como por \code{W6\_G03}. Esto no es tanto una contradicción como una
tensión antes no señalada entre dos nociones distintas de un ``buen'' corredor: las
pruebas de mérito (y la intuición del lector) favorecen a \code{W6\_G02} porque atiende
la mayor demanda absoluta no satisfecha a una escala que parece una línea real de
transporte; el objetivo W5 favorece a \code{W6\_G03} porque su conjunto pequeño y
uniformemente de brecha alta de anclas le da la mejor tasa \emph{promedio} de
focalización ponderada por demanda. Una metodología que reporta una función
multiobjetivo formal no debería reportar selectivamente solo la lente de evaluación que
favorece su ejemplo narrativo preferido; esta asimetría se corrige aquí en lugar de en
una revisión futura.

Probar una variante de cobertura absoluta de $f_1$ (\cref{sec:res-w6}) resuelve esto de
una manera más útil que simplemente elegir una lente sobre la otra: no es que
\code{W6\_G03} sea en secreto el candidato más fuerte bajo una ponderación ``más
justa'', sino que la formulación de tasa de $f_1$ estaba estructuralmente sesgada hacia
candidatos pequeños y geográficamente compactos, ya que un corredor corto puede
permanecer enteramente dentro de un bolsón homogéneo de alta necesidad mientras que uno
más largo necesariamente promedia algo de territorio de paso de menor necesidad para
seguir siendo una trayectoria conectada y construible. El hecho de que \code{W6\_G03}
ganara bajo \emph{todas} las ponderaciones objetivo probadas (\cref{sec:res-w6}) era en
sí mismo la señal: una compensación multiobjetivo genuina no debería producir un
candidato que domine absolutamente sin importar los pesos. Corregir $f_1$ a un total
absoluto ponderado por demanda elimina el artefacto y devuelve un frente de Pareto
genuino de cuatro vías ---\code{W6\_G02} y \code{W6\_G01} atienden sustancialmente más
demanda total no satisfecha, \code{W6\_G03} sigue siendo el más eficiente en demanda por
kilómetro, y ninguno domina a los demás. La elección que en realidad presenta el
generador de esta tesis a un planificador es, por tanto, real: financiar menos
corredores, más grandes, que atiendan más necesidad absoluta no satisfecha, o un
conector pequeño y muy focalizado que es eficiente por kilómetro pero modesto en
escala. Recomendar a \code{W6\_G02} como resultado principal es defendible en estos
términos ---para una única inversión de capital insignia, el impacto absoluto es el
criterio más relevante para la decisión--- pero la tesis debe decirlo explícitamente en
lugar de dejar que una métrica de tasa no examinada implique que las dos lentes
simplemente medían lo mismo de forma distinta. Junto a ambos,
el generador también arroja conectores de menor eficiencia (\code{W6\_G00} y
\code{W6\_G01}) que atienden necesidad real pero puntúan bajo en demanda por
kilómetro. La limitación residual es arquitectónica ---el número fijo de anclas, el
agrupamiento espacial y el tronco del árbol de expansión no pueden alcanzar cada
bolsón de brecha alta--- y ahora se mide en lugar de solo afirmarse. Una segunda
limitación independiente, también ahora medida en lugar de solo afirmada, es que el
$f_1$ basado en tasa de la función objetivo necesita el complemento de cobertura
absoluta anterior para leerse correctamente; la metodología de producción sigue
reportando por defecto la versión de tasa (\cref{sec:w5}), de modo que cualquier
re-ejecución futura de esta evaluación debería calcular y reportar ambas.

\section{La reconstrucción es un proxy débil del mérito de un corredor}
\label{sec:disc-reconstruction}

La ilustración más clara de la distinción Pregunta-A/Pregunta-B es el fracaso del
generador en reconstruir el ferrocarril construido. En las retro-pruebas enmascaradas
no re-descubre las líneas premium: enmascarar todo el nivel premium da
\SI{15}{\percent} de traslape medio de trazado, y enmascarar líneas férreas
individuales da esencialmente cero (\cref{sec:res-w8}); el corredor de la Línea~4 se
empareja con apenas \num{0.05} de exhaustividad. Tomado ingenuamente, esto parece un
fracaso. Leído correctamente, se acerca al resultado esperado y es solo evidencia débil
en contra del generador.

La razón es mecánica además de epistémica. Enmascarar una sola línea férrea apenas
mueve la accesibilidad, porque las rutas de autobús paralelas permanecen y la frontera
servido/no servido apenas se desplaza; no se forma una brecha nueva fuerte en el
corredor excluido, de modo que el generador de frontera no tiene dónde anclar. Este es
el mismo mecanismo que hace degenerar la retro-prueba de tronco de demanda en la Toluca
solo de autobuses (\cref{sec:transf-validation}). Epistémicamente, incluso una
puntuación de reconstrucción perfecta no establecería el mérito de un corredor, porque
las líneas construidas codifican restricciones ajenas a la demanda. La comparación en
las ciudades de transferencia es la imagen especular instructiva: donde la red
existente ya cubre bien la demanda, el generador se traslapa fuertemente con ella
(Toluca \SI{75}{\percent}, Aguascalientes \SI{54}{\percent}) ---coincidencia por
preferencia revelada, no cobertura nueva. El encuadre honesto, por tanto, es que el
valor del generador reside en la Pregunta~B, donde está parcialmente validado, no en la
Pregunta~A, donde la coincidencia o el desacuerdo son solo débilmente diagnósticos.

\section{Implicaciones de equidad}
\label{sec:disc-equity}

La equidad se trata como un término explícito y separable en lugar de como una
consecuencia implícita. La formulación aditiva (\cref{eq:final}) permite a la tesis
reportar la priorización con y sin el bono de equidad, y el análisis de sensibilidad
muestra que el ordenamiento es robusto al peso de equidad ($\alpha$): la estructura
general es fija, y $\alpha$ solo cambia cuáles \ageb{}s específicos encabezan la lista.
Sustantivamente, el diagnóstico y el término de equidad apuntan en la misma dirección
---las zonas de brecha alta son desproporcionadamente periferias densas y de alto
rezago social---, de modo que priorizar la demanda no atendida y priorizar la
desventaja están, en \zmg, en gran medida alineados y no en tensión. La corrección del
signo y la imputación del índice de marginación (documentada en el capítulo de datos)
importa aquí: un término de equidad que hubiera premiado silenciosamente a las zonas
menos marginadas habría invertido esta conclusión. Esta alineación no es meramente una
coincidencia dentro de \zmg de la geografía demográfica de una sola ciudad: la
\cref{tab:equity-transfer} muestra la misma correlación positiva entre brecha de
cobertura y puntuación de equidad, de magnitud similar, en ambas ciudades de
transferencia bajo su proxy de equidad ordinal más grueso, lo cual es cierta evidencia
de que la alineación refleja una regularidad más general de la demanda de transporte no
atendida y la desventaja social coubicándose en periferias mexicanas en crecimiento, y
no un artefacto específico de los propios indicadores de \zmg.

\section{Limitaciones}
\label{sec:disc-limitations}

Varias limitaciones acotan los resultados más allá de las ya señaladas. La superficie
de demanda descansa en un parámetro de decaimiento transferido y una impedancia
euclidiana; un modelo gravitacional de distancia de red calibrado con datos de encuesta
locales tensaría las magnitudes absolutas. La instantánea \textsc{gtfs} es un único
punto en el tiempo, por lo que la capa de oferta no capta la variación de frecuencia de
servicio a lo largo del día. La compuerta de factibilidad del generador es sensible a
la densidad de paradas del \textsc{gtfs} de origen, que en la auditoría se manifiesta
como una cuenta de factibles artificialmente baja para redes concesionadas de paradas
muy próximas. Lo más importante: los corredores \emph{generados} se validan solo de
forma indirecta ---por necesidad, no redundancia y eficiencia de demanda--- y no contra
afluencia observada, porque no existe afluencia contrafactual para líneas que nunca se
construyeron. La capa de diagnóstico, en cambio, tiene un anclaje fuera de muestra en la
Línea~4, aunque ese anclaje es en sí mismo un caso único reportado sin intervalo de
confianza (\cref{sec:res-w8}) y debe leerse como corroboración direccional y no como una
prueba estadísticamente cuantificada. Esta asimetría es el límite honesto de la
contribución: un diagnóstico
validado, y un generador caracterizado cuya efectividad se mide en sus propios términos.

Dos limitaciones adicionales, surgidas durante una revisión del 2026-09-09, se
documentan aquí. Primero, la afirmación de ``sin fuga'' del reentrenamiento supervisado
(\cref{sec:w3}) se había verificado solo contra la exclusión directa de columnas de
oferta de transporte; una ablación de características y una reevaluación bloqueada
espacialmente de seguimiento (\cref{sec:w3}) confirmaron que la preocupación era real y
no hipotética: remover las características de población cuesta \num{0.165} puntos de
PR-AUC de prueba, y una retención espacial a nivel municipal cuesta \num{0.247} puntos
respecto a la configuración de división aleatoria de producción. El PR-AUC de
producción de \num{0.877}--\num{0.883} debe leerse en consecuencia como una cota
superior bajo condiciones de evaluación optimistas. Segundo, un intento durante esta
revisión de reproducir la comparación de $R^2$ entre el prior y el calibrado de W2 sacó
a la luz una no-reproducibilidad que una investigación de seguimiento acotó pero no
cerró por completo. Re-ejecutar el mismo código de calibración contra la base de datos
en vivo devolvió $\beta^\star = \num{1.0025}$ y $n=2140$ pares de zonas, en lugar del
$\beta^\star = \num{1.2005}$, $n=\num{1993}$ reportado en toda esta tesis y en tres
ejecuciones históricas previas de calibración registradas en
\code{features.w2\_calibration}. La investigación de seguimiento descartó varias causas
candidatas en lugar de confirmar una: (i) las producciones agregadas, la demanda de
transporte y la tasa vehicular media de \code{features.ageb\_trip\_ends} coinciden
exactamente con los valores reportados en \cref{sec:res-w1}; (ii) el modelo
gravitacional W1 (\cref{sec:w1}) se re-ejecutó desde cero y reprodujo el mismo conteo de
filas de la matriz OD y el mismo flujo total (neto del corte documentado de
almacenamiento disperso $\ge\num{0.5}$, que por diseño excluye una larga cola de flujos
muy pequeños y no es en sí mismo un defecto); (iii) \code{src/w2\_gravity\_calibration.py}
no se ha modificado desde antes de que se produjera el resultado de \num{1.2005}; y sin
embargo recalibrar contra la salida de W1 recién confirmada sigue devolviendo
$\beta^\star=\num{1.0025}$, no \num{1.2005}. Con la matriz OD y las capas de generación
de viajes ahora confirmadas de forma independiente como reproducibles y el script de
calibración sin cambios, el candidato restante es la ruta de datos del lado de la
encuesta \textsc{eod} (la unión espacial \ageb-a-zona o las geometrías de zona
subyacentes) a un grano más fino de lo que puede revelar el \emph{conteo de
coincidencias} \ageb-a-zona (que es estable en 1,822/63); no sobrevive ninguna
instantánea de la base de datos en vivo de la fecha de la ejecución original de
\num{1.2005} contra la cual comparar directamente, razón por la cual esto pudo acotarse
pero no cerrarse. Dada la magnitud de re-derivar cada etapa posterior (W1, W3, W4, W6,
W7, W8) bajo un $\beta$ distinto, la cifra actualmente reproducible de \num{1.0025} no
se adoptó en esta tesis; se conserva la cifra reportada de \num{1.2005}, consistente con
tres ejecuciones históricas previas, en espera de recuperar el entorno computacional
original o de una decisión explícita de re-derivar los resultados de la tesis bajo la
nueva cifra. Esto se señala en lugar de resolverse en silencio porque incide
precisamente en el tipo de afirmación de reproducibilidad que la \cref{sec:datos-integridad}
por lo demás se esmera en documentar.
